% !TEX root = sub-main.tex
\documentclass[main.tex]{subfiles}

\begin{document}
    \begin{frame}
        \frametitle{What Can Explain JP Morgan's XVA Loss?}
        
        \begin{figure}
            \centering
            \resizebox{\textwidth}{!}{%
            \begin{tikzpicture}
                \colorlet{funding-color}{wtf-blue}
\colorlet{rate-color}{wtf-red}

\tikzmath{
    \arrowyoffset = 7;
    \arrowcenteryoffset = 7;
    \toparrowyoffset = \arrowcenteryoffset + \arrowyoffset;
    \bottomarrowyoffset = \arrowcenteryoffset - \arrowyoffset;
    \arrowtoboxpadding = 5;
}
\coordinate (center) at (0,0);
\node (dealer) at (center) [
    draw,
    fill=white!80!gray,
    very thick,
    minimum width=3cm,
    minimum height=2cm,
    align=center,
    rounded corners=.50cm,
] {JP Morgan};

\node (counterparty) at ($(dealer.west) + (-5, 0)$) [
    draw,
    very thick,
    minimum width=3cm,
    minimum height=2cm,
    align=center
] {Counterparty};

\node (exchange) at ($(dealer.east) + (5, 0)$) [
    draw,
    dashed,
    very thick,
    minimum width=3cm,
    minimum height=2cm,
    align=center
] {Exchange};

\node (funding) at ($(dealer.south) + (0, -4)$) [
    draw,
    very thick,
    minimum width=3cm,
    minimum height=2cm,
    align=center
] {Funding desk};

% Dealer - Counterparty relations
\draw[->, thick] 
    ([yshift=\bottomarrowyoffset pt, xshift=-\arrowtoboxpadding pt]dealer.south west) -- 
    ([yshift=\bottomarrowyoffset pt, xshift= \arrowtoboxpadding pt]counterparty.south east)
    node[pos = 0.35, fill=white, inner sep=2pt] {\footnotesize
        floating leg
    };
\draw[<-, thick] 
    ([yshift=\toparrowyoffset pt, xshift=-\arrowtoboxpadding pt]dealer.south west) -- 
    ([yshift=\toparrowyoffset pt, xshift= \arrowtoboxpadding pt]counterparty.south east)
    node[pos = 0.65, fill=white, inner sep=2pt] {\footnotesize
        fixed leg
    };

% Dealer - Exchange relations
\draw[<-, thick] 
    ([yshift=\bottomarrowyoffset, xshift= \arrowtoboxpadding pt]dealer.south east) -- 
    ([yshift=\bottomarrowyoffset, xshift=-\arrowtoboxpadding pt]exchange.south west)
    node[pos = 0.65, fill=white, inner sep=2pt] {\footnotesize
        floating leg
    };
\draw[->, thick] 
    ([yshift=\toparrowyoffset, xshift= \arrowtoboxpadding pt]dealer.south east) -- 
    ([yshift=\toparrowyoffset, xshift=-\arrowtoboxpadding pt]exchange.south west)
    node[pos = 0.35, fill=white, inner sep=2pt] {\footnotesize
        fixed leg
    };

\draw[->, thick, draw=funding-color] 
    ([yshift=-\bottomarrowyoffset, xshift= \arrowtoboxpadding pt]dealer.north east) -- 
    ([yshift=-\bottomarrowyoffset, xshift=-\arrowtoboxpadding pt]exchange.north west)
    node [pos=0.4, fill=white, inner sep=0pt] {
        \includegraphics[width=0.5cm]{\graphicsfolder/source-of-funding-costs/collateral}
    };
\draw[<-, thick, draw=rate-color] 
    ([yshift=-\toparrowyoffset, xshift= \arrowtoboxpadding pt]dealer.north east) -- 
    ([yshift=-\toparrowyoffset, xshift=-\arrowtoboxpadding pt]exchange.north west)
    node [pos=0.6, fill=white, inner xsep=2pt, inner ysep=0] {
        \footnotesize OIS
    };

% Dealer - Funding desk relations
\draw[<-, thick, draw=funding-color]
    ([xshift=-\arrowyoffset, yshift=-\arrowtoboxpadding]dealer.south) --
    ([xshift=-\arrowyoffset, yshift= \arrowtoboxpadding]funding.north)
    node[midway, fill=white, inner xsep=2pt, inner ysep=0, align=center, rotate=90] {
        \footnotesize Funding%
    };
\draw[->, thick, draw=rate-color]
    ([xshift=\arrowyoffset, yshift=-\arrowtoboxpadding]dealer.south) --
    ([xshift=\arrowyoffset, yshift= \arrowtoboxpadding]funding.north)
    node[midway, fill=white, inner xsep=2pt, inner ysep=0, align=center, rotate=90] {
        \footnotesize OIS + S%
    };
            \end{tikzpicture}
            }   
        \end{figure}

        \begin{tikzpicture}
            [
                remember picture,
                overlay,
                shift={([yshift=3.3cm,xshift=3cm]current page.south west)},
                align=left
            ]
            \node[font=\small]{
                \textbf{The rate effect} \\
                
                floating leg $\downarrow$

                \\[0.2cm]

                \hspace{0.8em} \rotatebox[origin=c]{180}{$\Lsh$}
                \hspace{0.4em} 
                \raisebox{-.45em}{\includegraphics[width=0.4cm]{\graphicsfolder/source-of-funding-costs/collateral}}
                $\uparrow$

                \\[0.2cm]

                \hspace{1.6em} \rotatebox[origin=c]{180}{$\Lsh$}
                \hspace{0.4em} funding $\uparrow$
            };
        \end{tikzpicture}

        \begin{tikzpicture}
            [
                remember picture,
                overlay,
                shift={([yshift=3.3cm,xshift=-3cm]current page.south east)},
                align=left
            ]
            \node[font=\small]{
                \textbf{The funding effect} \\
                
                default intensity $\uparrow$

                \\[0.2cm]

                \hspace{0.8em} \rotatebox[origin=c]{180}{$\Lsh$}
                \hspace{0.4em} S $\uparrow$

                \\[0.2cm]

                \hspace{1.6em} \rotatebox[origin=c]{180}{$\Lsh$}
                \hspace{0.4em} funding rate $\uparrow$
            };
        \end{tikzpicture}

        \note<1>[item]{
            Risk.net speculates that a contributing factor to the loss 
            may have been uncollateralised swap agreements.
        }
        \note<1>[item]{
            We explained in the paper how asymmetrical collateral agreements 
            can be a source of funding costs,
            and the corresponding graphic can explain the hypothesis of Risk.net.
        }
        \note<1>[item]{
            Not all corporate clients has the ability 
            to offer the large amount of collateral that banks need,
            and therefore the trades with those clients would be uncollateralized.
            Since corporate clients are mostly hedging loan exposures,
            banks are often biased towards having more uncollateralized pay-floating trades.
            These swaps are hedged with other financial institutions, 
            who require collateralization.
        }
        \note<1>[item]{
            Central banks decreasing rates will decrease the value of floating legs. 
            This increases the value of JP Morgan's trade with its corporate clients.
            It decreases the mark-to-market of the hedges however;
            therefore, JP Morgan will receive calls for collateral.
            If the portfolio had a somewhat equal amount of payer- and receiver-swaps,
            the effect would be close to zero.
        }
        \note<2>[item]{
            Since creditworthiness is decreasing, 
            JP Morgan will have to fund itself at increasing rates.
            Evidently, having these uncollateralized pay-floating swaps,
            is a source of WWR as the mark-to-market of the hedges is negatively correlated
            with the funding rate. 
            The opposite of that is shown in our swap example, 
            where the large upfront price were due in a period
            where our imaginary firm has a low funding spread. 
        }
        \note<2>[item]{
            This shows a very common way of funding costs appearing in derivatives portfolios.
            Another source that we covered thoroughly could appear
            if for example the payment dates of the hedging swap differed from the actual swap,
            such that the cash flows from market movements were unaligned.
        }
    \end{frame}
\end{document}
